\section{De ReAct a Plan-and-Solve}

\subsection{Limitaciones de ReAct}

ReAct procesa problemas complejos mediante el ciclo Thought $\rightarrow$ Action $\rightarrow$ Observation, pero tiene un problema clave:

\begin{warningbox}{Context Issues}
Cuando el ciclo de ReAct es especialmente largo, aparecen problemas graves de contexto:
\begin{itemize}
    \item El modelo de lenguaje sufre \textbf{Lost in the Middle} — se pierde en un contexto extenso y largo
    \item Olvida cuál era la tarea original del usuario
    \item Olvida qué acciones ha realizado y qué resultados obtuvo
    \item No logra producir un Action razonable para el siguiente paso
\end{itemize}
\end{warningbox}

\subsection{Plan-and-Solve Prompting}

Plan-and-Solve es una mejora del CoT Zero-shot clásico. Su forma de Prompting es la siguiente:

\begin{quote}
\textit{``Let's first understand the problem and devise a step-by-step plan to solve it. Then let's carry out the plan and solve the problem step by step.''}
\end{quote}

En comparación con el clásico ``Let's think step by step'', Plan-and-Solve introduce explícitamente el concepto de Plan dentro del Prompting, lo que produce una mejora de unos cuantos puntos porcentuales en problemas complejos.

\subsection{Dynamic Plan-and-Solve}

La comunidad hizo mejoras adicionales a Plan-and-Solve e introdujo el \textbf{carácter dinámico}:

\begin{importantbox}{Plan mental vs. interacción real}
El Plan se forma en la mente (un plan mental), sin interacción real con el entorno. Al ejecutarse de verdad se genera nuevo feedback, y en ese momento se necesita \textbf{Re-Plan} — incorporar el nuevo feedback al mecanismo de planificación para hacer una corrección dinámica.
\end{importantbox}

El flujo completo de Dynamic Plan-and-Solve:
\begin{enumerate}
    \item \textbf{Plan}: el Planner forma, de arriba hacia abajo, un plan (List of Steps / To-Do List)
    \item \textbf{Execute}: cada step se entrega a un ReAct Agent para que lo ejecute
    \item \textbf{Re-Plan}: después de la ejecución, con base en el nuevo feedback, se decide si terminar directamente o volver a planificar
\end{enumerate}

La entrada de Re-Plan incluye: el problema original + el Plan original + los pasos ejecutados previamente y sus resultados.

\begin{knowledgebox}{Consideración económica: combinación de modelos fuertes y débiles}
Desde el punto de vista económico, se puede usar un \textbf{modelo fuerte} como Planner y para el Re-Plan (porque la planificación determina el límite superior de calidad del conjunto), y un \textbf{modelo débil} para la ejecución concreta con ReAct (porque los pasos concretos de ejecución son relativamente simples). Así se logra un equilibrio entre costo y efecto.
\end{knowledgebox}

\subsection{Resumen de la sección}

De ReAct a Plan-and-Solve y luego a Dynamic Plan-and-Solve se aprecia una línea de evolución clara: de la interacción puramente reactiva, a la introducción explícita de la planificación, y luego al ajuste dinámico de la planificación para cerrar el gap entre el «plan mental» y la interacción real.


